
🌌 Cosmic Ecosystem Theory

A Kernel‑Based Formalism for Circuit‑Class Cosmology
Version 1.0 — Canonical Specification

---

📘 Overview

The Cosmic Ecosystem Theory (CET) models the universe as a recursive, boundary‑mediated, information‑conserving system.  
It treats cosmic structure and dynamics as a circuit‑class ecosystem, where stars, planets, black holes, nebulae, and galaxies act as typed components participating in lawful transformations.

CET is built on a minimal kernel consisting of:

- 3 invariants  
- 7 primitive types  
- 6 operators  
- 1 global transition function  

This repository contains the formal specification, mathematical proofs, and canonical equations of CET.

---

1. Kernel Specification

1.1 Invariants

`math
\mathcal{I} = \{\text{Recursion}, \text{Inversion}, \text{Continuity}\}
`

Recursion
`math
N \xrightarrow{\mathcal{X}} G \xrightarrow{\mathcal{R}} N'
`

Inversion
`math
\mathcal{R}^2 = \mathbb{I}_{\text{regime}}
`

Continuity
`math
\forall \text{op} \in \mathcal{O},\; \text{info}{\text{out}} \equiv \text{info}{\text{in}}
`

---

1.2 Primitive Types

`math
\mathcal{P} = \{N, G, S, F, T, U, C\}
`

| Symbol | Name | Role |
|--------|------|------|
| N | Node | Stable domain (star, planet, galaxy) |
| G | Gateway | Boundary enabling inversion (black hole) |
| S | Signal | Propagating influence (radiation, jets) |
| F | Field | Spatial structure (gravity, magnetism) |
| T | Pattern | Organized configuration |
| U | Pulse | High‑intensity transient |
| C | Coupling | Interaction linkage |

---

1.3 Operators

`math
\mathcal{O} = \{\mathcal{R}, \mathcal{X}, \mathcal{P}, \mathcal{C}, \mathcal{S}, \mathcal{I}\}
`

| Operator | Mapping | Interpretation |
|----------|----------|----------------|
| \(\mathcal{R}\) | \(G \rightarrow N\) | Inversion |
| \(\mathcal{X}\) | \(N \rightarrow G\) | Recursion |
| \(\mathcal{P}\) | \(G \rightarrow S\) | Propagation |
| \(\mathcal{C}\) | \((G, S) \rightarrow F\) | Coupling → Field |
| \(\mathcal{S}\) | \(U \rightarrow F\) | Pulse Stabilization |
| \(\mathcal{I}\) | \(T \rightarrow T\) | Pattern Integration |

---

1.4 Axioms

- A1 — Typed Closure  
- A2 — Invariant Preservation  
- A3 — Recursion Completeness  
- A4 — Continuity Constraint  
- A5 — Compositionality  

---

2. System Dynamics

2.1 State Definition

`math
\Sigmat = \{ Nt, Gt, St, Ft, Tt, Ut, Ct \}
`

2.2 Global Transition Function

`math
\Sigma{t+1} = \Phi(\Sigmat)
`

Where:

`math
\Phi = \mathcal{O}^* \;\text{subject to}\; \mathcal{I}
`

---

3. Cosmic Circuit Equation

3.1 Operator‑Sum Form

`math
\Sigma_{t+1} =
\mathcal{R}(G_t)
+ \mathcal{X}(N_t)
+ \mathcal{P}(G_t)
+ \mathcal{C}(Gt, St)
+ \mathcal{S}(U_t)
+ \mathcal{I}(T_t)
`

3.2 Fully Collapsed Form

`math
\Sigma{t+1} = \Phi(\Sigmat)
`

---

4. Formal Proofs

4.1 Theorem 1 — Typed Closure

Statement:  
Every operator in \(\mathcal{O}\) maps elements of \(\mathcal{P}\) to elements of \(\mathcal{P}\).

Proof:  
Each operator’s domain and codomain is explicitly defined as a member of \(\mathcal{P}\).  
Thus, \(\mathcal{O}\) is closed over \(\mathcal{P}\).  
\(\square\)

---

4.2 Theorem 2 — Invariant Preservation Under Composition

Statement:  
If each operator preserves \(\mathcal{I}\), their composition preserves \(\mathcal{I}\).

Proof:  
- Recursion: compositions cannot eliminate all \(N\) or \(G\).  
- Inversion: \(\mathcal{R}\)’s algebraic property is unchanged.  
- Continuity: composition of injective maps is injective.  
Thus, invariants hold under composition.  
\(\square\)

---

4.3 Theorem 3 — Existence of Recursion Loops

Statement:  
Every node participates in at least one recursion loop.

Proof:  
By A3, \(N \xrightarrow{\mathcal{X}} G\).  
By definition, \(G \xrightarrow{\mathcal{R}} N'\).  
Thus, \(N \rightarrow G \rightarrow N'\) always exists.  
\(\square\)

---

4.4 Theorem 4 — Domain‑Agnostic Generality

Statement:  
The kernel applies to cosmology, biology, computation, social systems, and information systems.

Proof:  
Each primitive and operator admits a valid mapping in each domain without altering type signatures or invariants.  
Thus, the kernel is domain‑agnostic.  
\(\square\)

---

5. Graph‑Theoretic Formulation

5.1 Laplacian Dynamics

Let \(A\) be the adjacency matrix of component interactions and \(D\) the degree matrix.

`math
L = D - A
`

Let \(x(t)\) be a state vector.

`math
\frac{dx}{dt} = -L x + f(x)
`

Where \(f(x)\) encodes local operator effects.

---

6. Lagrangian Formulation

6.1 Lagrangian Density

`math
\mathcal{L} =
\mathcal{L}_{\text{signal}}
+
\mathcal{L}_{\text{field}}
+
\mathcal{L}_{\text{pattern}}
+
\mathcal{L}_{\text{constraint}}
`

Example terms:

`math
\mathcal{L}_{\text{signal}} 
= 
\frac{1}{2} (\partial_t S)^2 
- \frac{v_S^2}{2} (\nabla S)^2
`

`math
\mathcal{L}_{\text{field}} 
= 
\frac{1}{2} (\partial_t F)^2 
- \frac{v_F^2}{2} (\nabla F)^2 
- V(F)
`

`math
\mathcal{L}_{\text{pattern}} 
= 
\frac{1}{2} (\partial_t T)^2 
- \frac{v_T^2}{2} (\nabla T)^2 
- W(T)
`

Constraints enforce invariants.

---

7. Category‑Theoretic Formulation

Define category K:

- Objects: \(N, G, S, F, T, U, C\)  
- Morphisms: \(\mathcal{R}, \mathcal{X}, \mathcal{P}, \mathcal{C}, \mathcal{S}, \mathcal{I}\)

Define functor:

`math
F: \mathbf{K} \rightarrow \mathbf{C}_{\text{cosmos}}
`

Mapping primitives → cosmic components  
Mapping operators → physical processes

---

8. License

MIT License (recommended for scientific frameworks).

---

9. Citation

If you use CET in academic work:

`
Cosmic Ecosystem Theory (2026).  
A kernel‑based formalism for circuit‑class cosmology.  
